\subsubsection{Sparse Diffusion Targeting}

A key observation in our early experiments is that directly applying diffusion enhancement to the entire hard mask leads to overuse of generative modification, which damages surrounding background consistency and introduces seam artifacts.
Therefore, we explicitly restrict diffusion to a sparse subset of difficult pixels instead of the whole candidate region.

Our pipeline first applies ProPainter as the main restoration module and then constructs a \texttt{diffusion\_target} mask for local enhancement.
The target mask is obtained from candidate regions after excluding borrowable areas, followed by confidence gating, morphological filtering, and connected-component cleanup.
When the target region becomes too large, we enforce an upper bound on the ratio between diffusion-target pixels and hard-mask pixels, so that only the most difficult pixels are retained for diffusion refinement.
The difficulty score combines low support fraction and high color inconsistency, thereby prioritizing pixels that are less reliably restored by propagation-based methods.

This sparse design is supported by quantitative evidence.
Here, \texttt{target/hard} denotes the ratio between the number of pixels selected for diffusion refinement and the number of pixels in the hard mask.
We report both the mean ratio over all evaluated samples and the 90th percentile (p90), which reflects the upper-end behavior in high-coverage cases.
Before introducing the ratio constraint, the average \texttt{target/hard} ratio was approximately \(0.915\), indicating that diffusion was effectively applied to almost the entire hard region.
After the sparse strategy was introduced, this ratio was reduced to \(0.35\) or \(0.25\), confirming that diffusion was restricted to a much smaller and more selective subset of pixels.

\begin{table}[t]
\centering
\caption{Effect of sparse targeting on diffusion coverage. Lower values indicate more selective diffusion refinement.}
\label{tab:sparse_target}
\begin{tabular}{lcc}
\toprule
Setting & target/hard mean & target/hard p90 \\
\midrule
Before sparsification & 0.915 & 0.994 \\
Sparse target (\(0.35\)) & 0.350 & 0.350 \\
Sparse target (\(0.25\)) & 0.250 & 0.250 \\
\bottomrule
\end{tabular}
\end{table}

Among these settings, the sparse-\(0.35\) configuration was considered the most balanced in terms of quality and stability, and was therefore kept as the default setting inside the diffusion-enhanced branch.

\begin{figure}[t]
    \centering
    \includegraphics[width=\linewidth]{figures/part3_sparse_target_vis.png}
    \caption{Sparse diffusion targeting. From left to right: hard mask, diffusion target, and enhanced result.}
    \label{fig:sparse_target}
\end{figure}

\subsubsection{Improved Boundary Blending}

Besides overuse of diffusion, another major artifact in our early results is the visible halo or bright boundary around refined regions.
Our analysis suggests that this issue comes from the interaction between VAE decoding artifacts and an overly rigid alpha blending scheme, which effectively pastes boundary noise back into the restored frame.

To alleviate this problem, we adopted a distance-transform-based alpha blending strategy.
Instead of using a hard or poorly shaped transition band, the new design assigns low alpha values near the outer boundary and gradually increases the replacement weight toward the interior of the refined region.
This creates a much smoother geometric transition between the enhanced patch and the original restored frame.

In practice, this modification substantially reduces halo artifacts and improves single-frame visual naturalness.
It is one of the clearest technical improvements obtained in Part 3.
Nevertheless, this improvement mainly benefits spatial quality on individual frames, and does not by itself resolve the temporal instability introduced by frame-wise diffusion enhancement.

\begin{figure}[t]
    \centering
    \includegraphics[width=\linewidth]{figures/part3_boundary_blending.png}
    \caption{Improved boundary blending. The distance-transform alpha reduces visible halo and seam artifacts.}
    \label{fig:boundary_blending}
\end{figure}

\subsubsection{Failure Analysis of Flow-based Temporal Propagation}

Although localized diffusion can improve difficult regions on key frames, it also introduces temporal inconsistency because enhanced frames may become noticeably sharper than their neighbors.
To reduce this flickering effect, we further explored flow-based residual propagation, where high-quality details generated on key frames are propagated to nearby frames using optical flow.

However, this direction did not become our final solution.
Both our experiments and the underlying task structure suggest that flow-based propagation is fundamentally weak for object removal and background reconstruction.
The reason is that this task is not merely about transporting visible pixels across time.
Instead, it often requires generating previously occluded background content that is never visible in adjacent frames.
Optical flow is effective for estimating correspondence between already visible appearances, but it is not designed to hallucinate missing content.

This mismatch becomes most severe near occlusion boundaries.
After object removal, the missing regions often involve disocclusion, one-to-many correspondence, and invalid backward mappings, making the inverse problem ill-posed.
Moreover, 2D flow models are easily corrupted by parallax, non-rigid motion, motion blur, and interpolation around high-frequency boundaries, which explains why seam and outside-region metrics are particularly sensitive to propagation strength.
In addition, propagation is recursive by nature, so even small alignment errors accumulate over time and manifest as temporal flicker.

Our quantitative results are consistent with this theoretical analysis.

\begin{table*}[t]
\centering
\small
\setlength{\tabcolsep}{4pt}
\caption{Effect of propagation strength on the gain--risk trade-off. The best score is achieved in a weak-propagation regime.}
\label{tab:flow_failure}
\begin{tabular}{lccccccc}
\toprule
Method & Radius & Blend & Outside $\downarrow$ & Seam $\downarrow$ & Leakage $\downarrow$ & Hard Gain $\uparrow$ & Score $\uparrow$ \\
\midrule
\texttt{sdxl-flow-on}          & 1 & 0.18  & 0.156413 & 3.535168 & 0.040955 & 0.000000 & -3.746414 \\
\texttt{freeinpaint-prev-on}   & 1 & 0.18  & 0.002539 & 0.144994 & 0.000000 & 0.078447 & 0.009360 \\
\texttt{freeinpaint-tuned-off} & 0 & 0.04  & 0.001069 & 0.086774 & 0.000000 & 0.053647 & 0.019451 \\
\texttt{c2-off}                & 0 & 0.03  & 0.002072 & 0.125207 & 0.000000 & 0.062710 & -0.001858 \\
\texttt{c3-off}                & 0 & 0.05  & 0.000977 & 0.051433 & 0.000000 & 0.021996 & -0.008418 \\
\texttt{c4-off}                & 0 & 0.035 & 0.000988 & 0.059160 & 0.000205 & 0.000000 & -0.060148 \\
\texttt{c5-off}                & 0 & 0.04  & 0.001675 & 0.111595 & 0.001482 & 0.048658 & -0.015955 \\
\bottomrule
\end{tabular}
\end{table*}

Strong propagation settings increased risk terms substantially, while the best-performing setting appeared in a weak-propagation regime rather than a strong-propagation regime.
For example, compared with a stronger flow-on configuration, a more conservative weak-propagation setting reduced \texttt{outside} from \(0.002539\) to \(0.001069\) and \texttt{seam} from \(0.144994\) to \(0.086774\), while improving the composite score from \(0.009360\) to \(0.019451\).
This supports the conclusion that optical-flow propagation should only be treated as a weak auxiliary mechanism, rather than the primary solution for temporal stabilization in this task.

More importantly, the comparison across different settings suggests that the dominant contradiction of this task is not controlled by a single propagation hyperparameter.
Even when propagation is weakened or disabled, the trade-off among \texttt{hard\_gain}, \texttt{seam}, and \texttt{leakage} remains difficult to resolve.
This is consistent with the underlying nature of object removal: the task simultaneously requires visible-content borrowing, invisible-content generation, and strong boundary control.
Pure flow propagation addresses the first requirement well, but remains weak on the latter two.

Overall, our experiments indicate that the diffusion-enhanced branch can improve some difficult single-frame cases, especially after sparse targeting and improved alpha blending.
However, it does not consistently outperform the Part 2 baseline in overall evaluation, mainly because temporal stability remains the dominant bottleneck.
Therefore, we keep the Part 2 pipeline as the default system and regard the diffusion branch as an exploratory extension with informative failure cases.